\scanpage{96}
% archive foliation: 97
\uline{Th} Soit $G=G(n)$ groupe de BT tronqué d'échelon $n$, $M$ un module \unsure{qu. coh.} annulé par $p^n$.
\figure{\begin{tikzpicture}\draw (0,0)--(0,0.7);\end{tikzpicture}}
On a
\diagram{\left\{\begin{array}{l}
\underline{\operatorname{Ext}}^1_{\Lambda_n}(G,M)=0\\
\underline{\operatorname{Ext}}^2_{\Lambda_n}(G,M)\xrightarrow{\sim}\underline{\operatorname{Ext}}^2_{\mathbf Z}(G,M)
\end{array}\right.}

\uline{NB} $\underline{\operatorname{Ext}}^2_{\mathbf Z}(G,M)=0$ si $2$ inv., pour $\forall G$ fini loc. libre; cela simplifie la théorie des déformations pour les BT en car. $p\ne2$ !).
